\chapter{系统实现}
\label{chap:implementation}

本章详细介绍基于驱动程序替换的固件仿真系统的具体实现。系统采用 Python 作为主要开发语言，结合 CodeQL 静态分析引擎、LangGraph 智能体框架和 Unicorn 仿真引擎，实现了从静态分析到动态测试反馈的完整工具链。

\section{实现概述}
\label{sec:implementation_overview}

\subsection{技术栈}

本系统的实现基于以下技术栈：

\textbf{（1）静态分析技术}

\begin{itemize}
    \item \textbf{CodeQL CLI 2.15+}：GitHub 提供的代码分析平台，支持 C/C++、Java、Python 等多种语言的语义分析
    \item \textbf{Python 3.10+}：作为主要开发语言，实现分析脚本和工具集成
\end{itemize}

\textbf{（2）大语言模型与智能体}

\begin{itemize}
    \item \textbf{LangGraph 0.2+}：构建状态ful 的图式工作流，支持循环和条件分支
    \item \textbf{LangChain 0.1+}：LLM 应用开发框架，提供结构化输出和工具调用
    \item \textbf{OpenAI API / Anthropic API}：调用 GPT-4o/Claude-3-Opus 等大语言模型
\end{itemize}

\textbf{（3）仿真与测试}

\begin{itemize}
    \item \textbf{Unicorn 2.0.2+}：轻量级多架构 CPU 模拟器框架
    \item \textbf{AFL++ 4.08c+}：覆盖率引导的模糊测试框架
    \item \textbf{ARM GCC 10.3+}：交叉编译工具链
\end{itemize}

\textbf{（4）构建与依赖管理}

\begin{itemize}
    \item \textbf{CMake 3.21+}：跨平台构建系统
    \item \textbf{Poetry / pip}：Python 依赖管理
\end{itemize}

\subsection{目录结构}

系统采用模块化的目录结构，便于维护和扩展：

\begin{lstlisting}[style=CStyle,abovecaptionskip=0pt,belowcaptionskip=0pt]
lcmhalmcp/
├── src/
│   ├── collector/           # 静态分析模块
│   │   ├── codeql/        # CodeQL 查询规则
│   │   ├── models.py      # 数据模型定义
│   │   └── analyzer.py    # 分析器实现
│   ├── builder/            # 替换与编译模块
│   │   ├── agents/        # LLM 智能体
│   │   ├── prompts/       # 提示词模板
│   │   └── compiler.py    # 编译集成
│   ├── feedback/           # 动态反馈模块
│   │   ├── emulator.py     # 仿真器封装
│   │   ├── lcmhal.py       # LCMHal 接口实现
│   │   └── agents/        # 反馈智能体
│   └── shared/            # 共享组件
│       ├── cache.py       # 缓存管理
│       └── utils.py       # 工具函数
├── tests/                  # 测试用例
├── configs/                # 配置文件
└── docs/                   # 文档
\end{lstlisting}

\subsection{关键依赖}

系统的核心依赖包括：

\begin{lstlisting}[style=CStyle,abovecaptionskip=0pt,belowcaptionskip=0pt]
# pyproject.toml
[tool.poetry.dependencies]
python = "^3.10"
codeql = "^2.15"
langgraph = "^0.2"
langchain = "^0.1"
langchain-openai = "^0.1"
anthropic = "^0.18"
unicorn = "^2.0"
pydantic = "^2.0"
\end{lstlisting}

\section{驱动函数接口静态分析与识别模块实现}
\label{sec:collector_impl}

本节详细介绍静态分析模块的具体实现，包括 CodeQL 查询规则、数据模型和分析流程。

\subsection{CodeQL 查询实现}

\subsubsection{MMIO 实例定位查询}

针对第三章介绍的三种 MMIO 定位模式，本节给出相应的 CodeQL 查询实现。

\textbf{查询一：直接地址强转模式}

\begin{lstlisting}[style=CStyle,abovecaptionskip=0pt,belowcaptionskip=0pt]
/**
 * 查询所有直接地址强转的 MMIO 实例
 * 模式：#define USART1 ((USART_TypeDef *)0x40011000)
 */
import cpp

from Variable v, Type t, Expr init
where
  // 类型名以 _TypeDef 或 _Type 结尾
  t.getName().regexpMatch(".*_TypeDef|.*_Type") and
  // 变量类型为指向该结构体的指针
  v.getType().(PointerType).getBaseType() = t and
  // 获取初始化表达式
  init = v.getInitializer() and
  // 初始化涉及整数常量
  init.getType() instanceof IntegralType and
  // 地址在常见的外设范围内
  init.getValue().toInt() in [0x40000000 .. 0x5FFFFFFF]
select v, t, init, "Direct cast: " + v.getName()
\end{lstlisting}

\textbf{查询二：间接初始化模式}

\begin{lstlisting}[style=CStyle,abovecaptionskip=0pt,belowcaptionskip=0pt]
/**
 * 查询通过句柄间接初始化的 MMIO 实例
 * 模式：huart1.Instance = USART1;
 */
import cpp

from Assignment assign, Variable instanceVar, Variable halVar
where
  // 赋值语句
  assign = instanceVar.getAnAssignedValue() and
  // 目标变量是结构体指针类型
  instanceVar.getType() instanceof PointerType and
  instanceVar.getType().(PointerType).getBaseType().getName()
      .regexpMatch(".*HandleType") and
  // 源变量是 MMIO 结构体实例
  halVar.getName().regexpMatch(".*_BASE|USART.*|GPIO.*") and
  assign.getRValue() = halVar.getAnAccess()
select assign, instanceVar, halVar,
       "Indirect init: " + instanceVar.getName() + " = " + halVar.getName()
\end{lstlisting}

\textbf{查询三：常量指针参数模式}

\begin{lstlisting}[style=CStyle,abovecaptionskip=0pt,belowcaptionskip=0pt]
/**
 * 查询以 MMIO 结构体指针为参数的函数
 * 模式：void HAL_UART_Transmit(USART_TypeDef *huart, ...)
 */
import cpp

from Function f, Parameter p
where
  // 函数名以 HAL_ 或库前缀开头
  f.getName().regexpMatch("HAL_.*|LL_.*|.*_Transmit|.*_Receive") and
  // 参数是指针类型
  p.getType() instanceof PointerType and
  p.getType().(PointerType).getBaseType().getName()
      .regexpMatch(".*_TypeDef|.*_Type") and
  // 参数属于该函数
  p.getFunction() = f
select f, p, "Function with MMIO param: " + f.getName()
\end{lstlisting}

\subsubsection{驱动函数定位查询}

\begin{lstlisting}[style=CStyle,abovecaptionskip=0pt,belowcaptionskip=0pt]
/**
 * 查询所有包含 MMIO 字段访问的函数
 */
import cpp

from Function f, FieldAccess fa
where
  // 字段访问
  fa.getTarget().getName().regexpMatch("DR|SR|CR.*|ISR|ICR|BRR") and
  // 字段访问所在的函数
  fa.getEnclosingFunction() = f and
  // 排除内联函数和宏定义
  not f.isInline() and
  not f.getName().matches("macro_%")
select f, fa.getAFunction(),
       "Driver function: " + f.getName() +
       " accesses " + fa.getTarget().getName()
\end{lstlisting}

\subsubsection{数据流追踪查询}

\begin{lstlisting}[style=CStyle,abovecaptionskip=0pt,belowcaptionskip=0pt]
/**
 * MMIO 数据流追踪配置
 * 从寄存器访问追踪到函数参数
 */
import cpp, semmle.code.cpp.dataflow.TaintTracking

class MMIOAccessConfig extends TaintTracking::Configuration {
  MMIOAccessConfig() { this = "MMIOAccessConfig" }

  // 定义源点：MMIO 结构体字段访问
  override predicate isSource(DataFlow::Node source) {
    exists(FieldAccess fa |
      // 字段名匹配典型的寄存器名称
      fa.getTarget().getName().regexpMatch("DR|SR|CR.*|ISR|ICR") and
      source.asExpr() = fa
    )
  }

  // 定义汇点：函数参数
  override predicate isSink(DataFlow::Node sink) {
    exists(Parameter p |
      p.getFunction().getName().regexpMatch("HAL_.*|LL_.*") and
      sink.asParameter() = p
    )
  }

  // 定义额外的污点步骤：通过指针解引用传播
  override predicate isAdditionalTaintStep(
    DataFlow::Node pred, DataFlow::Node succ
  ) {
    exists(PointerDereferenceExpr pde |
      succ.asExpr() = pde and
      pred.asExpr() = pde.getOperand()
    )
  }
}

// 执行查询
from MMIOAccessConfig config, DataFlow::PathNode source, DataFlow::PathNode sink
where
  config.hasFlowPath(source, sink)
select source, sink, "Data flow: register access -> parameter"
\end{lstlisting}

\subsection{数据模型实现}

\subsubsection{函数信息模型}

\begin{lstlisting}[style=Python,abovecaptionskip=0pt,belowcaptionskip=0pt]
# src/collector/models.py
from dataclasses import dataclass
from typing import List, Optional
from enum import Enum

class FunctionCategory(str, Enum):
    """驱动函数分类"""
    RECV = "RECV"          # 接收数据
    SEND = "SEND"          # 发送数据
    IRQ = "IRQ"            # 中断处理
    INIT = "INIT"          # 初始化
    CTRL = "CTRL"          # 控制
    POLL = "POLL"          # 轮询
    DMA = "DMA"            # DMA 操作
    NODRIVER = "NODRIVER"  # 非驱动函数

@dataclass
class Parameter:
    """函数参数"""
    name: str
    type: str
    is_pointer: bool

@dataclass
class FunctionInfo:
    """函数基本信息"""
    name: str
    file_path: str
    line_start: int
    line_end: int
    signature: str
    return_type: str
    parameters: List[Parameter]
    category: Optional[FunctionCategory] = None
    source_code: Optional[str] = None
    caller_functions: List[str] = None
    callee_functions: List[str] = None

    def __post_init__(self):
        if self.caller_functions is None:
            self.caller_functions = []
        if self.callee_functions is None:
            self.callee_functions = []
\end{lstlisting}

\subsubsection{MMIO 访问点模型}

\begin{lstlisting}[style=Python,abovecaptionskip=0pt,belowcaptionskip=0pt]
# src/collector/models.py

from enum import Enum

class AccessType(str, Enum):
    """访问类型"""
    READ = "read"
    WRITE = "write"
    BOTH = "both"

@dataclass
class MMIOAccess:
    """MMIO 访问点"""
    struct_type: str        # 如 USART_TypeDef
    field_name: str         # 如 DR, SR
    access_type: AccessType
    line_number: int
    enclosing_function: str
    context_code: str       # 上下文代码片段（前后5行）

    @property
    def access_summary(self) -> str:
        """生成访问摘要"""
        return f"{self.struct_type}->{self.field_name} ({self.access_type})"
\end{lstlisting}

\subsubsection{项目分析结果模型}

\begin{lstlisting}[style=Python,abovecaptionskip=0pt,belowcaptionskip=0pt]
# src/collector/models.py
from datetime import datetime

@dataclass
class ProjectAnalysisResult:
    """项目分析结果"""
    project_path: str
    codeql_db_path: str
    functions: List[FunctionInfo]
    mmio_accesses: List[MMIOAccess]
    call_graph: dict
    driver_functions: List[str]
    analysis_timestamp: str

    @classmethod
    def from_dict(cls, data: dict) -> 'ProjectAnalysisResult':
        """从字典创建实例"""
        # 反序列化逻辑
        pass

    def to_dict(self) -> dict:
        """转换为字典"""
        # 序列化逻辑
        pass

    def save(self, filepath: str):
        """保存到文件"""
        import json
        with open(filepath, 'w', encoding='utf-8') as f:
            json.dump(self.to_dict(), f, indent=2, ensure_ascii=False)

    @classmethod
    def load(cls, filepath: str) -> 'ProjectAnalysisResult':
        """从文件加载"""
        import json
        with open(filepath, 'r', encoding='utf-8') as f:
            return cls.from_dict(json.load(f))
\end{lstlisting}

\subsection{分析流程实现}

\subsubsection{数据库构建}

\begin{lstlisting}[style=Python,abovecaptionskip=0pt,belowcaptionskip=0pt]
# src/collector/analyzer.py
import subprocess
from pathlib import Path

class CodeQLAnalyzer:
    def __init__(self, project_path: str):
        self.project_path = Path(project_path)
        self.db_path = self.project_path / ".codeql" / "database"

    def create_database(self, language: str = "cpp") -> bool:
        """
        创建 CodeQL 数据库

        Args:
            language: 目标语言（cpp, java, python等）

        Returns:
            是否成功
        """
        cmd = [
            "codeql", "database", "create",
            str(self.db_path),
            "--language=" + language,
            "--source-root=" + str(self.project_path),
            "--overwrite"
        ]

        try:
            result = subprocess.run(
                cmd,
                cwd=self.project_path,
                capture_output=True,
                text=True,
                timeout=300  # 5分钟超时
            )
            return result.returncode == 0
        except subprocess.TimeoutExpired:
            print("Database creation timeout")
            return False
        except Exception as e:
            print(f"Database creation failed: {e}")
            return False
\end{lstlisting}

\subsubsection{查询执行}

\begin{lstlisting}[style=Python,abovecaptionskip=0pt,belowcaptionskip=0pt]
# src/collector/analyzer.py

class CodeQLAnalyzer:
    # ... 其他方法 ...

    def run_query(self, query_file: str) -> List[dict]:
        """
        执行 CodeQL 查询

        Args:
            query_file: 查询文件路径（.ql文件）

        Returns:
            查询结果列表
        """
        cmd = [
            "codeql", "query", "run",
            "--database=" + str(self.db_path),
            "--output=" + str(self.project_path / "results.bqrs"),
            query_file
        ]

        # 执行查询
        result = subprocess.run(cmd, capture_output=True, text=True)

        if result.returncode != 0:
            print(f"Query execution failed: {result.stderr}")
            return []

        # 转换为 CSV 格式
        csv_path = self.project_path / "results.csv"
        cmd_bqrs_decode = [
            "codeql", "bqrs", "decode",
            "--format=csv",
            "--output=" + str(csv_path),
            str(self.project_path / "results.bqrs")
        ]

        subprocess.run(cmd_bqrs_decode, capture_output=True)

        # 解析 CSV 结果
        import csv
        results = []
        with open(csv_path, 'r') as f:
            reader = csv.DictReader(f)
            for row in reader:
                results.append(row)

        return results
\end{lstlisting}

\section{本章小结}
\label{sec:implementation_summary}

本章详细介绍了基于驱动程序替换的固件仿真系统的具体实现。系统采用模块化设计，各组件通过明确的接口进行交互，实现了从静态分析到动态反馈的完整闭环。

主要实现包括：

\textbf{（1）静态分析模块实现}：基于 CodeQL 实现了 MMIO 实例定位、驱动函数识别和数据流追踪的查询规则，并设计了统一的数据模型用于结果存储和传递。

\textbf{（2）LLM 智能体实现}：基于 LangGraph 实现了具有状态管理能力的智能体工作流，通过精心设计的提示词引导 LLM 进行函数分类和代码生成。

后续章节将详细介绍其余模块的实现细节，包括替换策略的具体代码实现、编译集成的自动化流程、以及动态测试反馈闭环的仿真平台和反馈机制。
